% Context from chapters-tex/inverse-problems.tex; for mathematical reference, not compilation.
for some $\rho>0$. This is equivalent to the coefficient bound
\eql{\label{eq-source-cond-quad}\tag{$S_{\be,\rho}$}
	\sum_m \si_m^{-2\be} |\dotp{f_0}{v_m}|^2 \leq \rho^2 < +\infty.
}


The assumptions $\beta>0$ and $f_0\in\ker(\Phi)^\bot$ are part of the source condition; the coefficient bound alone leaves kernel components unconstrained. The bound resembles the Sobolev constraint in~\ref{eq-discr-sobol}. For example, consider a low-pass convolution $\Phi f = f \star h$ where the kernel $h$ has a polynomially decaying multiplier: $|\hat h_m| \sim 1/m^{\al}$ for large $m$. Since the vectors $v_m$ are Fourier modes, \eqref{eq-source-cond-quad} gives, up to constants and the Fourier normalization, a bound of the form
\eq{
	\sum_m \norm{m}^{2\al\be} |\hat f_m|^2 \leq \rho^2 < +\infty.
}
This is a Sobolev-type coefficient bound of smoothness order $\al\be$.


    
\paragraph{Sublinear convergence speed.}

The following theorem gives the convergence rate implied by this source condition.
 
With the noise bound $\norm{w} \leq \de$, recovery depends on the parameters $(\de,\rho,\be)$. Assuming $f_0 \in \ker(\Phi)^\bot$, we study the convergence of $f_\la$ to $f_0$ for data $y=\Phi f_0+w$ as $\de \rightarrow 0$. The analysis also determines how $\la$ should depend on $\de$.


\begin{thm}\label{thm-sublin-quad}
	Let $\delta>0$. Under the source condition~\eqref{eq-source-cond-quad} with $0 < \be \leq 2$, the solution of~\eqref{eq-regul-normsq} for $\norm{w} \leq \de$ satisfies
	\eq{
		\norm{f_\la - f_0} \leq C \rho^{\frac{1}{\be+1}} \de^{\frac{\be}{\be+1}}
	}
	for a constant $C$ which depends only on $\be$, and
	for a choice 
	\eq{
		\la \sim \de^{\frac{2}{\be+1}} \rho^{-\frac{2}{\be+1}}. 
	} 
\end{thm}

\begin{proof}
	The source condition implies $f_0\in\ker(\Phi)^\bot$.
	 
	We decompose
	\eq{
		f_\la = \Phi_\la^+(\Phi f_0 + w) = f_\la^0 + \Phi_\la^+ w 
		\qwhereq		
		f_{\la}^0 \eqdef \Phi_\la^+ ( \Phi f_0 ), 
	}
	Thus $f_\la = f_{\la}^0 + \Phi_\la^+ w$, and for any regularized inverse of the form~\eqref{eq-regul-pinv}
	\eql{\label{eq-rate-tikhon}
		\norm{f_\la - f_0} \leq \norm{ f_\la - f_{\la}^0 } + \norm{ f_{\la}^0 - f_0 } .
	}
    The noise-propagation term $\norm{f_\lambda-f_\lambda^0}$ measures the effect of measurement noise. For Tikhonov regularization it decreases as $\la$ increases and is often called the variance term.
	 
    The bias term $\norm{ f_{\la}^0 - f_0 }$ is independent of the noise and measures the smoothing of $f_0$. For Tikhonov regularization it increases with $\la$. Choosing $\la$ therefore balances bias against noise amplification.
	 
	Assuming 
	\eq{
		\foralls \si\geq 0, \quad  |\mu_\la(\si)| \leq C_\la
	}
	the variance term is bounded as
	\eq{
		\norm{ f_\la - f_{\la}^0 }^2 =  \norm{ \Phi_\la^+ w }^2 = \sum_m \mu_\la(\si_m)^2 |w_m|^2 \leq C_\la^2 \norm{w}_\Hh^2.
	}
	Using $\frac{|f_{0,m}|^2}{\si_m^{2\be}} = |z_m|^2$, we bound the bias term as follows:
	\eql{\label{eq-proof-ip-quad-1}
		\norm{ f_{\la}^0 - f_0 }^2 = \sum_m ( 1 - \mu_\la(\si_m) \si_m)^2 |f_{0,m}|^2
		= \sum_m \pa{ (1 - \mu_\la(\si_m) \si_m) \si_m^{\be} }^2  \frac{|f_{0,m}|^2}{\si_m^{2\be}}
		\leq D_{\la,\be}^2 \rho^2
	}
	where we assumed
	\eql{\label{eq-proof-ip-quad-2}
		\foralls \si\geq 0, \quad \abs{ (1 - \mu_\la(\si) \si ) \si^{\be} } \leq D_{\la,\be}.
	}
	For $\be>2$, the supremum over all $\sigma\geq0$ is infinite. On the bounded spectrum of a fixed operator, the bias remains bounded, but the resulting rate saturates.
	Putting~\eqref{eq-proof-ip-quad-1} and~\eqref{eq-proof-ip-quad-2} together, one obtains
	\eql{\label{eq-bound-global-quad}
		\norm{f_\la - f_0} \leq C_\la \de + D_{\la,\be} \rho. 
	}
	In the case of the regularization~\eqref{eq-regul-normsq}, one has $\mu_\la(\si) = \frac{\si}{\si^2+\la}$, and
	thus $( 1 - \mu_\la(\si) \si )\si^{\be} = \frac{\la \si^\be}{\si^2+\la}$. 
	For $\be \leq 2$, one verifies (see Figure~\ref{fig-bound-regul-variance} and~\ref{fig-ip-bound-regul}) that
	\eq{
		C_\la = \frac{1}{2\sqrt{\la}}
		\qandq
		D_{\la,\be} = c_\be \la^{\frac{\be}{2}},  
	}
	for some constant $c_\be$.
	 
    Balancing the two terms in~\eqref{eq-bound-global-quad} gives the correct order; optimizing over $\la$ would improve only the constant. Set $\frac{\de}{\sqrt{\la}} = \la^{\frac{\be}{2}} \rho$, which gives $\la = (\de/\rho)^{\frac{2}{\be+1}}$. Then
	\eq{
		\norm{f_\la - f_0} = O( \de/\sqrt{\la} ) = O( \de (\de/\rho)^{-\frac{1}{\be+1}} ) 
		= O( \de^{\frac{\be}{\be+1}} \rho^{\frac{1}{\be+1}} ) .
	}
\end{proof}



\begin{figure}[tbp]
\centering
\BookDrawing{tikz/inverse-problems-variance-bound}
\caption{Left: capping reciprocal singular values. Right: the Tikhonov bound $\mu_{\la}(\si) \leq C_\la = \frac{1}{2\sqrt{\la}}$.}
\label{fig-bound-regul-variance}
\end{figure}


\begin{figure}[tbp]
\centering
\BookDrawing{tikz/inverse-problems-bias-bound}
\caption{Bounding $\la \frac{\si^\be}{\la+\si^2} \leq D_{\la,\be}$.}
\label{fig-ip-bound-regul}
\end{figure}



Larger source orders $\be \leq 2$ therefore give faster convergence as $\norm{w}$ tends to zero.
The rate in Theorem~\ref{thm-sublin-quad} saturates: choosing $\be > 2$ gives the same general worst-case rate as $\be=2$. The best rate obtainable in this way is
\eq{
	\norm{f_\la - f_0} = O( \rho^{\frac{1}{3}} \de^{\frac{2}{3}} ).
}
Alternative spectral filters $\mu_\la$, together with sufficiently large $\be$, can give the rate $\norm{f_\la - f_0} = O( \de^{1-\kappa} )$ for arbitrarily small $\kappa>0$.
 
The capped inverse in Figure~\ref{fig-bound-regul-variance}, left, avoids this finite-order saturation, whereas Tikhonov regularization is limited to source orders $\beta\leq2$. Quadratic regularization is easier to implement: its variational formulation leads to a linear system and does not require an SVD.
 
For compact operators with nonclosed range, no uniform linear rate follows over these source balls from finite-order source conditions alone. Linear rates are possible in finite dimension with a fixed smallest positive singular value, or under different structural assumptions, such as those used for sparse $\ell^1$ regularization in Chapter~\ref{chap-sparse-regul}.

                                                                
                                                                
                                                                
\section{Quadratic Regularization}

We now turn from infinite-dimensional recovery theory to numerical methods in finite dimension.

    
\paragraph{Convex regularization.}

As in~\eqref{eq-regul-inv}, we reconstruct a high-resolution image $f_0 \in \RR^N$ from noisy measurements $y = \Phi f_0 + w \in \RR^P$ by minimizing a convex regularized objective:
\eql{\label{eq-ip-regul}
	f_\la \in  \uargmin{ f \in \RR^N }\; \Ee(